\section{Authorship, development and AI assistance}
\label{sec:arch-authorship}

I developed BioSeqDB over approximately three years. Early work included
Python aligner prototypes written approximately two years ago. I
designed the system and implemented the BioSeqDB-specific algorithms,
representations and optimizations presented as my contributions in the
preceding chapters. Moving from those prototypes to a PostgreSQL extension
also required substantial integration, testing and performance work.

These contributions build on published algorithms and existing software.
The minimap2 mapping structure, prior SIMD minimizer work and adaptive block
alignment are important foundations.\autocite{minimap2_paper,
  groot_koerkamp_simd_minimizers_2025,block_aligner_paper}
The implementation includes ports and adaptations, as well as new work.
The basic columnar store comes from Citus and Hydra and retains code written
by their authors and maintainers, as described in
Section~\ref{sec:pg-storage-role}.\autocite{citus_columnar_storage,hydra_columnar_source}
PostgreSQL, SeqAn, Highway and oneTBB also provide substantial components.
The original contribution lies in the specific algorithms, changes and
integration identified in this thesis. It does not imply authorship of
the underlying published algorithms or of all source code in the repository.

\subsection{Where AI tools were used}
\label{sec:arch-ai-assistance}

I used AI coding tools under requirements and designs that I supplied.
Their work included:
\begin{itemize}
  \item \textbf{Source reorganization.} Long files were divided into modules,
    functions were renamed, and code was adjusted for Clang-based checks.
    Many present file and function names were produced during this work.
  \item \textbf{Development tools.} AI helped implement benchmarking utilities
    and code generators according to my designs, and assisted with
    implementation and optimization tasks.
  \item \textbf{Technical writing.} AI generated working documentation.
\end{itemize}

The large reorganizations were intended to preserve the algorithms while
changing names and file boundaries. They nevertheless introduced bugs.
In my assessment, diagnosing and correcting regressions from AI-assisted
work consumed approximately two months. Numerical code was particularly
sensitive: a seemingly equivalent change to a mask, integer width or
coordinate calculation could change results. A refactoring that compiles
is therefore insufficient to establish numerical correctness.

AI assistance accelerated parts of development despite these costs. The
research questions, project-specific designs and decisions about which
ideas to pursue were mine. I directed and corrected the assisted work and
retain responsibility for the final implementation and the claims made here.

Automatically generated documentation is not an authoritative specification.
Some descriptions were incorrect, and others became stale after later
changes. For this chapter, generated descriptions were checked against the
implementation, the cited technical sources and the stated tests. Generated
names and reports alone do not establish how a feature works.

\section{Validation limits of the complete system}
\label{sec:arch-validation-limits}

Validation of query-result caching, described in
Section~\ref{sec:bioseq-result-cache}, remains incomplete. It has targeted tests,
including comparisons with caching disabled, but these do not cover all
combinations of plans, settings and transaction histories. Incorrect reuse
and therefore incorrect query results remain a risk. The cache should be
treated as experimental.
